%!TEX root = main.tex
%!TEX spellcheck = en_US

\section{Introduction}
\label{sec:intro}

Ensuring the stability, security, and performance of network infrastructures
requires carefully configuring routers, switches, firewalls, etc. A misconfiguration on even a single network device can lead to black holes, unintended network access, inefficient routes, and a range of other issues~\cite{gill2011understanding, turner2010california}.

Although some configuration tasks are automated---especially in large networks~\cite{tian2019jinjing, ramanathan2023aura, mahimkar2022aurora, rasche2023esnet, gember-jacobson2015mpa}---many (changes to) configurations are manually written by human experts~\cite{abhashkumar2020aed}. Network engineers rely on extensive knowledge gained from years of experience with an organization's network(s), vast vendor documentation~\cite{ciscosupport, junosdocumentation}, industry best practices~\cite{manrs}, etc. To avoid problems, engineers must consider the interplay between (a change to) a configuration and related aspects of devices' configurations. This is challenging for humans, because even a single device's configuration is often hundreds or thousands of lines long with many dependencies between segments~\cite{benson2009complexitymetrics}.


Researchers and practitioners have developed numerous network configuration verifiers to address this challenge.
{\em Model checkers}~\cite{fogel2015general, beckett2017general, abhashkumar2020tiramisu, prabhu2020plankton, zhang2022sre, steffen2020netdice, ye2020hoyan, ritchey2000using,al2011configchecker, jeffrey2009model, batfish} analyze network-wide forwarding behaviors under various scenarios (\eg, link failures) to detect violations of network objectives (\eg, reachability). These tools rely on
%Tool developers must implement 
algorithms or formalisms that codify protocol- and vendor-specific configuration semantics as well as network-specific forwarding policies. These require significant manual effort to define---only some algorithms and policies can be automatically inferred~\cite{birkner2020config2spec, kheradmand2020anime, sharma2023prosper}---and are difficult to define correctly and completely~\cite{ye2020hoyan, birkner2021metha}; hence, model checkers have a high overhead to adoption.
{\em Consistency checkers}~\cite{kakarla2024diffy,kakarla2020finding, le2006minerals, feamster2005detecting, tang2021campion, le2008detecting, le2006characterization, batfish}, on the other hand, compare configurations across devices or against best-practices and flag deviations as potential issues.
% While this statistics-driven approach requires minimal manual effort, it lacks semantic depth.
While these statistics- or rule-driven approaches require minimal manual effort, they lack semantic depth and the ability to reason about intentional deviations.
For example, a consistency checker might flag a rarely used packet filter, even though the filter is required in some contexts to satisfy security objectives.


Large language models (LLMs)~\cite{deepseekai2024deepseekv3technicalreport, achiam2023gpt} offer a promising avenue to avoid the classic trade-off between semantic depth and ease-of-use in network verification. Pretrained on a vast, public corpus of data that includes protocol specifications and vendor documentation, LLMs develop a deep understanding of general network configuration principles. When presented with a specific, previously unseen device configuration at inference time, the LLM applies this broad knowledge to effectively translate the configuration into a reasoned judgment. This process hinges on the model's self-attention mechanism, which mathematically weighs the importance of every token in the configuration relative to every other token. By using attention, the model builds a rich, contextual map, understanding, for instance, that a specific access control list on line 50 directly impacts an interface definition on line 450. After creating this complex, attention-weighted internal representation, the model’s final layers distill it into a simple, probabilistic choice over a small set of output tokens, such as 'true' or 'false'. This ability to use attention to interpret complex dependencies and then collapse that understanding into a discrete assessment makes LLMs analogous to human experts applying broad domain knowledge to diagnose a specific problem.

However, recent efforts to leverage LLMs for configuration analysis~\cite{bogdanov2024leveraging,wang2024identifying,liu2024large, wang2024netconfeval, lian2023configuration} have missed an important insight:
human engineers do not examine entire configurations in a rigid, linear fashion; instead, engineers identify relevant configuration lines and dynamically ``jump'' to related parts of the configuration. Similarly, LLMs are more effective when provided with {\em concise, relevant context} that mirrors this human approach. Feeding an LLM the entire configuration at once risks exceeding token limits, and splitting configurations into arbitrary partitions~\cite{lian2023configuration} can obscure dependencies between different parts. Furthermore, overloading the model with irrelevant or excessive information can dilute its reasoning capabilities~\cite{lican,li2024long,qian2024long}, much like a human struggling to recall long-past details. We discuss these issues in more depth in Section~\ref{sec:background:context}.

The central challenge, therefore, is to bridge the gap between the reasoning capabilities of LLMs and the complex, interdependent nature of network configurations. An effective framework must enable an LLM to analyze configurations in a way that is both semantically deep and scalable, mirroring a human expert's ability to dynamically seek out relevant context without being overwhelmed by the sheer volume of a full configuration file.
In this paper, we propose the Context-Aware Prompting (\sysname{}) framework, a holistic solution that addresses this challenge through two integrated components:

\mypara{1. Automatic and Accurate Context Mining:}
To ground the LLM's reasoning, \sysname{} first performs automated context extraction. We model the configuration's nested structure as a hierarchical tree, where configuration stanzas and parameters (\eg, interface GE1/1/1) are intermediate nodes and terminal commands (\eg, mtu 9000) are leaf nodes. This structure enables the systematic identification of key context categories—neighboring, similar, and referenced configuration lines—ensuring a concise yet sufficient corpus is available for analysis. Furthermore, we distinguish pre-defined from user-defined values using an initial LLM query to classify the terminal value in a given line.

\mypara{2. Guided Interactive Prompting for Context Management:}
To prevent context overload and focus the model's attention, \sysname{} introduces an guided, interactive dialogue. While modern LLMs have large context windows, providing excessive or unfocused information is known to degrade reasoning quality, an issue distinct from hard token limits. Our framework mitigates this by having the LLM first explicitly request only the specific context it needs from the categories mined in the first component. This allows the model to incrementally build its understanding and enables more precise detection.


We evaluate \sysname{} on configurations from three diverse networks (Internet2 and two university campuses), structuring our analysis around four key questions. Our evaluation demonstrates that (1) \sysname{} achieves 100\% accuracy in detecting a wide range of known synthetic and real-world misconfigurations; (2) it can uncover novel, previously unknown issues in production networks, while also highlighting the challenge of false positives with some LLM backends; (3) the choice of LLM significantly impacts reasoning strategies and outcomes; and (4) \sysname{} substantially outperforms existing model checkers and consistency-based tools on a controlled set of complex errors.